\vspace{0.5em}\noindent\textbf{Sharing and Feedback}\par

This section summarizes my personal reflections after joining the
\textbf{First Cloud AI Journey (FCAJ)} program and completing the
\textbf{Internship Application Platform} project. The feedback focuses
on my learning experience, mentor and community support, program
strengths, and suggestions for future improvement.

\vspace{0.5em}\noindent\textbf{Overall Evaluation}\par

\vspace{0.5em}\noindent\textbf{1. Learning and Practice
Environment}\par

The program provided a good environment for learning Cloud and AI topics
through practical work. Instead of only studying theory, I had the
opportunity to build a project with multiple components, including
frontend, backend, realtime chat, AI service, databases, Docker,
Kubernetes, CI/CD, and observability.

I especially appreciated that the program encouraged learners to
research, implement, and record results independently. This approach
helped me understand more deeply how a system is designed, operated, and
deployed in a Cloud environment.

\vspace{0.5em}\noindent\textbf{2. Support from Mentors and Team
Admins}\par

Mentors and team admins provided helpful support during the learning
process. When I encountered issues, I could ask questions, receive
direction, and continue solving the problem by myself instead of only
receiving direct answers. This helped me improve debugging,
documentation reading, and independent problem-solving skills.

Information about the program, events, workshops, and reference
materials was also shared clearly, giving me suitable learning resources
for each stage.

\vspace{0.5em}\noindent\textbf{3. Alignment with My
Major}\par

The internship content was aligned with my Computer Science major and my
development direction. The project involved backend engineering,
databases, distributed systems, Cloud infrastructure, security, and AI
integration.

Through the project, I not only strengthened programming knowledge but
also learned more about non-functional requirements such as security,
scalability, reliability, monitoring, and cost control.

\vspace{0.5em}\noindent\textbf{4. Learning Opportunities and Skill
Development}\par

During the internship, I learned several important skills:

\begin{itemize}
\tightlist
\item
  Designing a multi-service system.
\item
  Building REST APIs with FastAPI.
\item
  Managing data with PostgreSQL, DynamoDB, Redis, and S3.
\item
  Integrating realtime chat with Socket.IO.
\item
  Integrating AI for CV and Job Description analysis.
\item
  Containerizing applications with Docker and deploying with Kubernetes.
\item
  Setting up CI/CD and observing systems with metrics, logs, and traces.
\item
  Writing bilingual technical reports with a clear structure.
\end{itemize}

These skills helped me gain a more practical view of building a
Cloud-native system instead of focusing only on individual features.

\vspace{0.5em}\noindent\textbf{5. Learning Culture and Community
Spirit}\par

The FCAJ community felt open and encouraged knowledge sharing. Online
events, workshops, and blog posts in the community gave me more
perspectives on AWS Cloud, Agentic AI, enterprise architecture, DevOps,
observability, and security.

I also valued the program's encouragement to share knowledge through
blogs and reports. This is a good way to turn learned knowledge into
content that can be presented and reviewed.

\vspace{0.5em}\noindent\textbf{Additional Questions}\par

\vspace{0.5em}\noindent\textbf{What was I most satisfied with during the
internship?}\par

I was most satisfied with the opportunity to build a relatively complete
project that covered business flows, AI integration, realtime
communication, deployment direction, and observability. The project
helped me understand the relationship between code, infrastructure, and
operations.

\vspace{0.5em}\noindent\textbf{What could the program improve for future
interns?}\par

The program could provide more checklists for each stage, such as
backend, frontend, AWS deployment, monitoring, and cost cleanup
checklists. This would help learners review their progress more easily
and avoid missing important parts.

In addition, sample project review sessions or discussions about common
AWS deployment issues would give learners more practical experience
before working on their own projects.

\vspace{0.5em}\noindent\textbf{Would I recommend the program to
friends?}\par

I would recommend the program to friends who are interested in Cloud,
backend, DevOps, or AI applications. FCAJ is suitable for learners who
want hands-on practice, product-oriented learning, and a better
understanding of how AWS services can work together in a real system.

\vspace{0.5em}\noindent\textbf{Suggestions and
Expectations}\par

\begin{itemize}
\tightlist
\item
  Provide project checklists for each milestone so learners can track
  progress more easily.
\item
  Add more examples about cost optimization and cleanup to avoid
  unnecessary AWS charges.
\item
  Share common issues when working with EKS, IAM, S3, CloudFront, RDS,
  and GitHub Actions.
\item
  Provide more guidance on technical report writing and architecture
  diagram presentation.
\item
  Organize short project review sessions so learners can receive
  feedback before finalizing reports.
\item
  Provide more career guidance for Cloud Engineer, Backend Engineer,
  DevOps Engineer, and AI Engineer paths.
\end{itemize}

\vspace{0.5em}\noindent\textbf{Career Impact}\par

After the program, I have a clearer direction toward Cloud/Backend
Engineering combined with AI integration. I realized that building a
modern application is not only about writing code. It also includes
architecture design, security, deployment, monitoring, cost management,
and operational documentation.

The program helped me become more confident in reading technical
documentation, analyzing requirements, selecting suitable services, and
presenting solutions in a structured way.

\vspace{0.5em}\noindent\textbf{Final Message}\par

I would like to thank the First Cloud AI Journey program, mentors, team
admins, and community for creating an environment where I could learn
and complete the project during the internship. The knowledge and
experience gained from this program are an important foundation for my
continued development in Cloud, backend, DevOps, and AI application
engineering.
